\section{模型的检验与灵敏度分析}

\subsection{硬约束与一致性检验}

由于真实最优任务位置、储能动作和用户满意度均不可观测，模型检验不采用准确率、召回率或混淆矩阵，而采用任务唯一性、时间可行性、资源残差、能源守恒、SOC 连续性和指标复算一致性。问题二的 GPU 容量、IT 功率、设施功率、网络时延和完成时限约束验证结果为 $1.0$，最大容量残差为 $0.0$；问题三的 SOC、充放电功率、效率、购售电边界及终端状态验证结果为 $1.0$；问题四的任务、储能和购售电综合约束验证结果同样为 $1.0$，最大容量残差为 $0.0$。因此，三个正式候选均通过声明范围内的完整硬约束复核。

预测部分采用固定起点递归检验，以避免验证区间和测试区间的未来真实值被读取。第 $2376$ 至 $2399$ 小时测试预测的区域—类型序列 MAE 均值为 $25.836515435967883$ 等效 GPU。按预声明 ACF 局部峰规则，可识别周期序列数为 $0.0$，因此周期性不能被作为已经得到数据支持的结论，季节滞后项仅作为预声明预测候选参与验证选择。

\subsection{碳权重灵敏度}

灵敏度实验以问题二固定调度的成本加碳代价为目标，基准目标值为 $-421673289.0742142$。当碳权重参数下降 $20\%$ 和 $10\%$ 时，目标变化率分别为 $-0.001790589840291007\%$ 和 $-0.0008952949201384358\%$；当碳权重上升 $10\%$ 和 $20\%$ 时，目标变化率分别为 $0.000895294920152571\%$ 和 $0.001790589840305142\%$。对应结果如图~\ref{fig:sens-carbon}所示。

\begin{figure}[htbp]
\centering
\includegraphics[width=0.88\textwidth]{sensitivity_carbon_weight.png}
\caption{碳权重变化下目标函数的灵敏度}
\label{fig:sens-carbon}
\end{figure}

图~\ref{fig:sens-carbon}显示，碳权重在给定扰动范围内变化时，目标变化率的绝对值均不超过结构化实验中记录的 $0.001790589840305142\%$。响应方向随参数增减保持一致，说明加权目标对碳权重变化具有稳定的单调响应。就目标数值而言，模型对碳权重较为稳健，但该结论只适用于当前固定调度和所给扰动范围，不能替代独立碳排放指标的复算。

\subsection{价格尺度灵敏度}

当价格尺度下降 $20\%$ 和 $10\%$ 时，目标变化率分别为 $20.001790589840287\%$ 和 $10.00089529492015\%$；当价格尺度上升 $10\%$ 和 $20\%$ 时，目标变化率分别为 $-10.00089529492015\%$ 和 $-20.001790589840287\%$。相应变化如图~\ref{fig:sens-price}所示。

\begin{figure}[htbp]
\centering
\includegraphics[width=0.88\textwidth]{sensitivity_price_scale.png}
\caption{价格尺度变化下目标函数的灵敏度}
\label{fig:sens-price}
\end{figure}

图~\ref{fig:sens-price}表明，价格尺度扰动几乎按相同量级传递至目标函数。与碳权重实验相比，价格尺度变化引起的目标变化率明显更大，因此模型对价格尺度敏感。由于基准目标为负值，变化率符号需要结合“购电支出减售电收入”的定义解释；负号变化并不等价于物理能耗下降，而表示统一财务目标的数值移动。

综合两组实验可得，目标函数对碳权重的局部扰动较稳健，对价格尺度变化则较敏感。实际应用中，碳权重可在给定范围内用于表达偏好而不显著改变固定调度目标，但购售电价格输入需要被重点校核，并应通过多价格情景重新运行能源策略，而不能只对最终成本作比例换算。灵敏度分析遵循参数扰动与输出复算的基本方法\cite{saltelli2008global}。
